\section{Week 8: Mode-2 HBZ and rANS Boundary}

Week~8 keeps the Week~7 product/replay freeze unchanged and moves only along
\claimid{OBL-SIG-HBZ-ENCODE-DECODE}.  The result is a compiled leaf inverse,
an exact focused/full extraction identity for the actual HBZ wrappers, a
compiled generic mode-2 arithmetic certificate, and a compiled concrete
actual-array table corollary.  The generated encoder/decoder loop refinement
remains the
gates to a full actual HBZ round trip, so the parent claim remains
\status{partial}.

\subsection*{Actual extraction and implementation graph}

The focused roots are
\procname{encode_hb_z1_prepare_jazz},
\procname{decode_hb_z1_apply_jazz},
\procname{rans_encode_jazz}, \procname{rans_decode_jazz}, and the two
mode-2 full wrappers.  Their generated closures form
\[
\begin{array}{c}
\procname{encode_hb_z1_mode2_full_jazz}\allowbreak
 \to \procname{_encode_hb_z1_full}\allowbreak
 \to \{\procname{_encode_hb_z1_prepare},
       \procname{_rans_encode},
       \procname{__copy_encoded_suffix}\},\\[2mm]
\procname{decode_hb_z1_mode2_full_jazz}\allowbreak
 \to \procname{_decode_hb_z1_full}\allowbreak
 \to \{\procname{_rans_decode},
       \procname{_decode_hb_z1_apply}\}.
\end{array}
\]
The wrappers fix $\mathsf{count}=1024$, $m=13$, and
$\mathsf{offset}=6$.  Regeneration checks show byte-identical generated
bodies for these procedures in the focused closure and in the Week~7 actual
signature pack/unpack closures.  The pinned source hashes include
\texttt{dff14e57\ldots3516cd} for \path{encoding.jinc},
\texttt{818bb254\ldots127b49} for \path{sparse_encoding.jinc}, and
\texttt{b1d4b2dd\ldots228429a} for \path{encoding_tables.jinc}.

\subsection*{Coefficient-to-symbol mapping}

For the first 1024 words, canonical mode-2 input means
\[
 -6 \leq \mathsf{signed}(hbz[i]) < 7.
\]
The actual prepare loop establishes
\[
 \mathsf{symbol}[i]=\mathsf{signed}(hbz[i])+6,
 \qquad 0\leq\mathsf{symbol}[i]<13,
\]
sets its bad word to zero, and preserves symbol bytes $[1024,2048)$.
The actual apply loop zero-extends the symbol and subtracts six modulo
$2^{32}$.  Its word-level inverse lemma handles negative input coefficients
and preserves coefficient words $[1024,2048)$.

The sequential actual-procedure theorem is
\path{encode_prepare_decode_apply_mode2_inverse}.  Its sole semantic input is
\path{canonical_hbz_mode2}; the conclusion contains prepare success, exact
first-1024 recovery, and both tail frames.  It is Hoare partial correctness,
not a termination claim.  The all-zero array supplies a compiled
precondition witness.

\subsection*{Concrete 13-symbol table certificate}

The concrete mode-2 intervals are:
\[
\begin{array}{c|rrrrrrrrrrrrr}
s&0&1&2&3&4&5&6&7&8&9&10&11&12\\ \hline
c_s&0&1&2&3&8&66&312&710&957&1016&1021&1022&1023\\
f_s&1&1&1&5&58&246&398&247&59&5&1&1&1
\end{array}
\]
Every frequency is positive and the sum is 1024.  The adjacent intervals
partition $[0,1024)$.  The authored certificate checks the generic encoder
field formulas, decoder packed start/frequency formulas, and the reciprocal
division arithmetic used by the mode-2 design.  The concrete actual-array
corollary is intentionally split out into
\theoryname{Mode2HbzSymbolWordsGenerated}; the external generator is only a
deterministic producer of EasyCrypt proof text, and the actual-array link is
accepted only because that generated theory itself fresh-compiles.

For normalized state $x$, the generic certificate also proves exactness of the
reciprocal multiply/high-half/shift computation.  Frequency-one symbols use
the implementation's $x-1$ quotient and compensated bias.  In every case the
actual fast expression equals
\[
 E_s(x)=\left\lfloor\frac{x}{f_s}\right\rfloor 1024
       +c_s+(x\bmod f_s),
\]
with a non-overflowing W64 product and a W32 state below $2^{31}$.

\subsection*{State machine, normalization, and inverse invariant}

The state lower bound is $L=2^{23}$ and the scale is $2^{10}$.  Encoding
traverses symbols backward.  Before $E_s$, bytes are emitted from the low end
of $x$ while $x\geq 2^{21}f_s$, using a cursor that moves backward.  The
final state is written as four little-endian bytes.  Decoding reads those
bytes, recovers symbols forward from the low 10-bit slot, applies
\[
 D_s(x')=\left\lfloor\frac{x'}{1024}\right\rfloor f_s
          +(x'\bmod1024)-c_s,
\]
and consumes normalization bytes while the state is below $L$.

The pure step theorem is now compiled.  In
\theoryname{Mode2RansCore}, \path{pure_rans_step_inverse} proves
$D_s(E_s(x))=x$, and \path{hbz_fast_step_decode_inverse} replaces
$E_s$ with the concrete reciprocal implementation formula under
$0\leq s<13$ and $1\leq x<\mathsf{hbz\_xmax}(s)$.
The remaining generated-loop invariant must additionally relate the
encoder's backward byte stack $[off,count)$ to the decoder's forward input,
preserve $L\leq x<2^{31}$ at symbol boundaries, and track failure before the
four-byte state area.  A prototype for the actual suffix copy discharged its
body invariant but did not finish the quantified loop-exit projection under
fresh compilation; no such theorem is retained as compiled evidence.

\subsection*{Success-conditioned boundary and failure branch}

Canonical coefficient bounds do not imply buffer success.  If the backward
encoder exhausts its fixed workspace, the actual full encoder returns
$\mathsf{size}=0$.  The target theorem must therefore have the form
\[
 \mathsf{size}=0\quad\lor\quad
 \left(4\leq\mathsf{size}\leq1024\;\land\;
       \mathsf{decode\_bad}=0\;\land\;
       decoded[0..1024)=original[0..1024)\right).
\]
Neither decoder success nor output equality may be assumed.  A successful
actual encoder/decoder witness is also required before the parent can be
promoted to \status{proved}.

The authored proof hashes are
\texttt{909e4eed\ldots0d865f} (prepare),
\texttt{03ac5595\ldots521d4} (apply),
\texttt{7da906bc\ldots92e55} (leaf composition),
\texttt{8bdecef8\ldots6ad58} (table arithmetic),
\texttt{3f664576\ldots{}d6ba90} (generated table certificate),
\texttt{6efb58b5\ldots5ba80} (actual extraction boundary), and
\texttt{a4a10f55\ldots{}a5838c} (pure rANS step).  Each is checked by
\texttt{easycrypt compile -script -no-eco} through
\path{scripts/verify-all.sh}; the extraction and generated certificate are
regenerated before those compilations.

\subsection*{Proof status, failed tactic, and security boundary}

\begin{center}
\begin{tabular}{l l}
\toprule
Obligation & Status \\
\midrule
HBZ prepare/apply and leaf frames & \status{proved} \\
Generic mode-2 table arithmetic certificate & \status{proved} \\
Concrete actual-array table certificate & \status{proved} \\
Pure arithmetic rANS step inverse & \status{proved} \\
Actual suffix copy and frame & \status{partial} \\
Actual encoder loop refinement & \status{partial} \\
Actual decoder loop refinement & \status{blocked} \\
Actual full HBZ left inverse & \status{partial} \\
\bottomrule
\end{tabular}
\end{center}

A monolithic \texttt{inline; sim} attempt is structurally inadequate: it
does not invent the backward-stack/forward-consumption relation, and it
obscures the size-zero branch.  Repeating larger full-procedure inlining is
rejected.  A weaker actual-encoder control prototype also stalled at its
nested normalization loop, after the outer post-loop commands and branch had
been separated with \texttt{wp}.  The next proof must refine each generated
loop against one explicit state/byte-stack model instead of adding another
wrapper.

An encode-to-decode left inverse is functional evidence only.  Canonical
parsing additionally needs malformed rejection or decode-to-re-encode
uniqueness.  Week~8 therefore proves neither a zero encoding delta nor
\claimid{OBL-SIGN-OUTPUT-TAIL-REACH}, full Verify correctness, signing
losslessness, or implementation EUF-CMA security.  Fresh verification uses
the aggregate script above and records target-specific logs under
\path{logs/}.
